% Estilo común de la presentación «Conteo de referencias wait-free: lo que dicen los experimentos».
% Lo cargan presentacion.tex (sin notas) y presentacion_con_notas.tex.
\usepackage[spanish,es-noshorthands,es-nodecimaldot,es-tabla]{babel}
\usepackage[T1]{fontenc}
\usepackage{tgheros}
\renewcommand{\familydefault}{\sfdefault}
\usepackage[scaled=0.9]{inconsolata}
\usepackage{booktabs,tabularx,array,graphicx,xcolor,etoolbox,amssymb,colortbl}
\usepackage{tikz}
\usetikzlibrary{calc,positioning,babel}
\hypersetup{pdftitle={Conteo de referencias wait-free: lo que dicen los experimentos},
  pdfauthor={},pdfcreator={},pdfproducer={},pdfsubject={Réplica · Euro-Par 2024},pdfkeywords={}}

% ---------- paleta ----------
\definecolor{azul}{HTML}{1F4E79}
\definecolor{azulclaro}{HTML}{9DB9D5}
\definecolor{tinta}{HTML}{1B2430}
\definecolor{gris}{HTML}{5B6573}
\definecolor{suave}{HTML}{EEF1F5}
\definecolor{ambar}{HTML}{B7791F}
\definecolor{ambarsuave}{HTML}{FBF6EC}
\definecolor{verde}{HTML}{1B7F3B}
\definecolor{rojo}{HTML}{B3261E}
% colores de los algoritmos (los mismos que en las gráficas)
\definecolor{cAutores}{HTML}{D62728}
\definecolor{cAnderson}{HTML}{9B4DCA}
\definecolor{cFolly}{HTML}{E69F00}
\definecolor{cWilliams}{HTML}{138A5E}
\definecolor{cLibs}{HTML}{3B8FD9}
\definecolor{cMutex}{HTML}{3A3A3A}
\newcommand{\alg}[2]{\textcolor{#1}{\rule{0.62em}{0.62em}}\,#2}

% ---------- plantilla ----------
\setbeamersize{text margin left=7mm,text margin right=7mm}
\setbeamertemplate{navigation symbols}{}
\setbeamercolor{normal text}{fg=tinta}
\setbeamercolor{structure}{fg=azul}
\setbeamercolor{alerted text}{fg=rojo}
\setbeamertemplate{itemize item}{\textcolor{azul}{\scriptsize$\blacksquare$}}
\setbeamertemplate{itemize subitem}{\textcolor{gris}{--}}
\setbeamertemplate{enumerate item}{\textcolor{azul}{\bfseries\insertenumlabel.}}
\setlength{\leftmargini}{1.2em}

% sección en curso (aparece encima del título) y fuente de cada diapositiva (al pie)
\gdef\laseccion{}
\gdef\lafuente{}
\newcommand{\fuente}[1]{\gdef\lafuente{#1}}
\BeforeBeginEnvironment{frame}{\gdef\lafuente{}}

\setbeamertemplate{frametitle}{%
  \vspace*{6pt}%
  \ifx\laseccion\empty\else{\color{azul}\fontsize{6.5}{7}\selectfont\bfseries\MakeUppercase{\laseccion}}\par\vspace{2pt}\fi
  {\color{tinta}\fontsize{12.5}{14}\selectfont\bfseries\insertframetitle}\par
  \vspace{3pt}{\color{azul!25}\rule{\linewidth}{0.5pt}}\vspace*{-2pt}}

\setbeamertemplate{footline}{%
  \leavevmode\hbox to\paperwidth{\hspace{7mm}%
    \begin{minipage}[b]{0.78\paperwidth}\raggedright\color{gris}\fontsize{5.6}{6.6}\selectfont
      \ifx\lafuente\empty\else Fuente: \lafuente\fi\end{minipage}%
    \hfill{\color{gris}\fontsize{5.6}{6.6}\selectfont Réplica · Euro-Par 2024\quad
      \textbf{\insertframenumber}\,/\,\inserttotalframenumber}\hspace{7mm}}\vspace{5pt}}

% ---------- bloques ----------
% \caja: idea clave (barra azul); \aviso: matiz o salvedad (barra ámbar)
\newcommand{\caja}[2][\linewidth]{%
  \begin{tikzpicture}\node[fill=suave,inner sep=5pt,text width=#1-10pt-2pt,align=justify,
      font=\small]
    (c){#2};\fill[azul] (c.north west) rectangle ([xshift=2pt]c.south west);\end{tikzpicture}}
\newcommand{\aviso}[2][\linewidth]{%
  \begin{tikzpicture}\node[fill=ambarsuave,inner sep=5pt,text width=#1-10pt-2pt,align=left,
      font=\footnotesize]
    (c){#2};\fill[ambar] (c.north west) rectangle ([xshift=2pt]c.south west);\end{tikzpicture}}
% \porque{TÍTULO}{texto}: recuadro «Por qué ocurre» con la teoría que explica un resultado
\newcommand{\porque}[3][\linewidth]{%
  \begin{tikzpicture}\node[draw=azul!20,line width=0.6pt,fill=white,inner sep=6pt,
      text width=#1-12pt-1.2pt,align=left,font=\footnotesize]
    {{\color{azul}\fontsize{6.5}{7.5}\selectfont\bfseries POR QUÉ OCURRE: \MakeUppercase{#2}}\par\smallskip #3};
  \end{tikzpicture}}
\newcommand{\sub}[1]{{\color{azul}\bfseries #1}\par\smallskip}
\newcommand{\cifra}[2]{\begin{minipage}[t]{\linewidth}\centering
  {\color{azul}\fontsize{22}{24}\selectfont\bfseries #1}\par\smallskip{\color{gris}\scriptsize #2}\end{minipage}}
\newcommand{\cmd}[1]{\colorbox{tinta}{\parbox{\dimexpr\linewidth-2\fboxsep}{\color{white}\ttfamily\scriptsize #1}}}
\newcommand{\pot}[1]{10\textsuperscript{#1}}
\newcommand{\por}{\texttimes}
% los dos repositorios de GitHub de la réplica (enlaces clicables)
\newcommand{\repo}[2]{\href{https://github.com/#1}{\ttfamily github.com/#1}\ #2}
\newcommand{\repos}{\begin{tabular}{@{}l@{}}
  \repo{dtamotu/replica-wait-free-refcount}{(rama \texttt{replica-completa})}\\
  \repo{davichiacustico-ux/replica-wait-free-refcount}{}\end{tabular}}

% veredictos
\newcommand{\etiqueta}[2]{{\setlength{\fboxsep}{1.4pt}\colorbox{#1}{\color{white}\fontsize{6}{7}\selectfont\bfseries #2}}}
\newcommand{\sirep}{\etiqueta{verde}{Se reproduce}}
\newcommand{\norep}{\etiqueta{rojo}{No se reproduce}}
\newcommand{\enparte}{\etiqueta{ambar}{En parte}}
\newcommand{\noconf}{\etiqueta{gris}{No se confirma}}

% ---------- portadas de parte ----------
\newcounter{parte}
\newcommand{\parte}[4]{%
  \stepcounter{parte}\gdef\laseccion{#1}%
  {\setbeamercolor{background canvas}{bg=azul}\setbeamertemplate{footline}{}%
  \frame[plain,c]{\hypertarget{parte\theparte}{}%
    {\color{azulclaro}\small\bfseries PARTE \theparte}\par\medskip
    {\color{white}\fontsize{22}{26}\selectfont\bfseries #2\par}\bigskip
    {\color{azulclaro}\normalsize #3\par}%
    \note{\input{notas/#4}}}}}
\newcommand{\notaparte}{}

% ---------- notas ----------
% \nota{sNN} carga notas/sNN.tex como nota del orador de la diapositiva.
\newcommand{\nota}[1]{\note{\input{notas/#1}}}
\setbeamertemplate{note page}{%
  \vspace*{6pt}%
  \begin{minipage}[t]{0.30\textwidth}\vspace{0pt}%
    \fcolorbox{azul!30}{white}{\insertslideintonotes{0.27}}\end{minipage}\hfill
  \begin{minipage}[t]{0.66\textwidth}\vspace{2pt}%
    {\color{azul}\fontsize{6.5}{7.5}\selectfont\bfseries NOTAS · DIAPOSITIVA \insertframenumber\,/\,\inserttotalframenumber}\par\smallskip
    {\color{tinta}\bfseries\small\insertframetitle}\end{minipage}\par\medskip
  {\color{azul!25}\rule{\textwidth}{0.5pt}}\par\smallskip
  \fontsize{8.6}{11}\selectfont\color{tinta}\setlength{\parskip}{4pt}\raggedright
  \insertnote}
